\documentclass[10pt,
% a4paper,
%twocolumn,
fleqn,
%landscape, 
% papersize,
dvipdfmx,
uplatex
]{jsarticle}



\def\maru#1{\textcircled{\scriptsize#1}}%丸囲み番号

% \RequirePackage[2020/09/30]{platexrelease}

%太字設定
\usepackage[deluxe]{otf}

\usepackage{emathEy}

\usepackage[g]{esvect}

%大きな文字
\usepackage{fix-cm}

%定理環境
\usepackage{emathThm}
%\theoremstyle{boxed}
\theorembodyfont{\normalfont}
\newtheorem{Question}{問題}[subsection]
\newtheorem{Q}{}[subsection]
\newtheorem{question}[Question]{}
\newtheorem{quuestion}{}[subsection]

%セクション，大問番号のデザイン
\renewcommand{\labelenumi}{(\arabic{enumi})}
\renewcommand{\theenumii}{\alph{enumii})}
\renewcommand{\thesection}{第\arabic{section}章}

%用紙サイズの詳細設定
\usepackage{bxpapersize}
\papersizesetup{size={80mm,45mm}}
\usepackage[top=0.7zw,bottom=0truemm,left=3truemm,right=133truemm]{geometry}
\usepackage[dvipdfmx]{graphicx}

%余白など
\usepackage{setspace} % 行間
\setlength{\mathindent}{1zw}
\setlength\parindent{0pt}


%色カラーに関する設定
\usepackage{color}
\definecolor{shiro}{rgb}{0.95703125,0.87109375,0.7421875}
\definecolor{kin}{rgb}{0.95703125,0.87109375,0.7421875}
\definecolor{orange}{rgb}{1,0.7,0.2}
\definecolor{bradorange}{rgb}{1,0.5,0}
\definecolor{pink}{rgb}{0.9176,0.5686,0.5960}
\definecolor{mizu}{rgb}{0.6156,0.8,0.9955}
\color{kin}
% \pagecolor{hukamido}

\usepackage{at}%図の配置
% \usepackage{wallpaper}

\begin{document}



\at(0cm,0cm){\includegraphics[width=8cm,bb=0 0 1920 1080]{./youtube/thumbnails/templates/smart_background/三角関数.jpeg}}
{\color{orange}\bf\boldmath\Large\underline{三角方程式の解の個数 Lv.2 }}\vspace{0.1zw}

\Large
\bf\boldmath 問．方程式

\Large
\vspace{0.3zw}
\hspace{0.0zw}
$\cos 2x+2a\sin x-4a-1=0\vspace{0.3zw}$

\Large
の$0\leqq x<2\pi$における\\
\hfill 異なる実数解の個数を求めよ．




\at(7.0cm,0.2cm){\small\color{bradorange}$\overset{\text{三角関数}}{\text{応用}}$}


\newpage



\at(0cm,0cm){\includegraphics[width=8cm,bb=0 0 1920 1080]{./youtube/thumbnails/templates/smart_background/三角関数.jpeg}}
{\color{orange}\bf\boldmath\Large\underline{三角方程式の解の個数 Lv.1 }}\vspace{0.1zw}

\Large
\bf\boldmath 問．方程式

\LARGE
\vspace{0.1zw}
\hspace{0.0zw}
$\sin ^2x+\cos x-\bunsuu{9}{4}+a=0\vspace{0.1zw}$

\Large
の$0\leqq x<2\pi$における\\
\hfill 異なる実数解の個数を求めよ．

\at(7.0cm,0.2cm){\small\color{bradorange}$\overset{\text{三角関数}}{\text{応用}}$}


\newpage



\at(0cm,0cm){\includegraphics[width=8cm,bb=0 0 1920 1080]{./youtube/thumbnails/templates/smart_background/三角関数.jpeg}}
{\color{orange}\bf\boldmath\LARGE\underline{半端な角の$\cos$の積}}\vspace{0.3zw}

\Large
\bf\boldmath 問．

\LARGE
\hspace{1zw}$\cos {20}^\circ \cos {40}^\circ \cos {80}^\circ$\vspace{0.5zw}

\Large
\hfill の値を求めよ．

\at(7.0cm,0.2cm){\small\color{bradorange}$\overset{\text{三角関数}}{\text{応用}}$}


\newpage



\at(0cm,0cm){\includegraphics[width=8cm,bb=0 0 1920 1080]{./youtube/thumbnails/templates/smart_background/三角関数.jpeg}}
{\color{orange}\bf\boldmath\LARGE\underline{半端な角の$\sin$の和}}\vspace{0.3zw}


\Large
\bf\boldmath 問．

\LARGE
\hspace{0.2zw}$\sin {10}^\circ +\sin {130}^\circ +\sin {250}^\circ$\vspace{0.5zw}

\Large
\hfill の値を求めよ．
\at(7.0cm,0.2cm){\small\color{bradorange}$\overset{\text{三角関数}}{\text{応用}}$}


\newpage



\at(0cm,0cm){\includegraphics[width=8cm,bb=0 0 1920 1080]{./youtube/thumbnails/templates/smart_background/三角関数.jpeg}}
{\color{orange}\bf\boldmath\LARGE\underline{$\sin 36^\circ$の値}}\vspace{0.3zw}

\normalsize
\bf\boldmath 問．$\theta ={36}^\circ$のとき，
$\sin 3\theta =\sin 2\theta$を示し，

\fontsize{45}{0} \selectfont
\vspace{0.3zw}
\hspace{1zw}$\sin {36}^\circ$
\vspace{0.2zw}

\normalsize
\hfill の値を求めよ．
\at(7cm,0.2cm){\small\color{bradorange}$\overset{\text{三角関数}}{\text{典型}}$}

\newpage

\at(0cm,0cm){\includegraphics[width=8cm,bb=0 0 1920 1080]{./youtube/thumbnails/templates/smart_background/三角関数.jpeg}}
{\color{orange}\bf\boldmath\LARGE\underline{和積公式の利用}}\vspace{0.3zw}\\
\LARGE 
\bf\boldmath 問．$0\leqq \theta \leqq \pi$のとき，次の方程式を求めよ．
\vspace{-0.3zw}
\[\hspace{-0.3zw} \sin 2\theta +\sin 3\theta +\sin 4\theta =0\]
\vspace{-0.3zw}
\at(7cm,0.2cm){\small\color{bradorange}$\overset{\text{三角関数}}{\text{典型}}$}

\newpage

\at(0cm,0cm){\includegraphics[width=8cm,bb=0 0 1920 1080]{./youtube/thumbnails/templates/smart_background/三角関数.jpeg}}
{\color{orange}\bf\boldmath\normalsize\underline{三角関数の最大値$〜\sin$と$\cos$の対称式$〜$}}\vspace{0.3zw}\\
\Large 
\bf\boldmath 問．次の関数の最大値と最小値を求めよ．
\vspace{-0.3zw}
\[\hspace{-0.3zw} y=2\sin {x}\cos {x}+\sin {x}+\cos {x}\]
\vspace{-2zw}
\[\hspace{8zw}\left(0\leqq x<2\pi \right)\]
\vspace{-0.3zw}
\at(7cm,0.2cm){\small\color{bradorange}$\overset{\text{三角関数}}{\text{典型}}$}

\newpage

\at(0cm,0cm){\includegraphics[width=8cm,bb=0 0 1920 1080]{./youtube/thumbnails/templates/smart_background/三角関数.jpeg}}
{\color{orange}\bf\boldmath\normalsize\underline{三角関数の最大$〜\sin$と$\cos$の$2$次式$〜$}}\vspace{0.3zw}\\
\Large 
\bf\boldmath 問．次の関数の最大値と最小値およびそのときの$x$の値を求めよ．
\vspace{-0.3zw}
\large
\[\hspace{-0.3zw}y=\sin ^2{x}+4\sin {x}\cos {x}+5\cos ^2{x}\]
\vspace{-2zw}
\[\hspace{10zw}\left(0\leqq x<2\pi\right)\]
\Large
\at(7cm,0.2cm){\small\color{bradorange}$\overset{\text{三角関数}}{\text{典型}}$}

\end{document}

